

\section{Introdução}
\label{sec:introducao}
O avanço dos modelos de linguagem (\textit{Large Language Models} - LLMs), transformou o paradigma de interação entre humanos e máquinas. No entanto, a ausência inicial de padrões de comunicação entre as LLMs e ferramentas externas gerou uma fragmentação técnica que motivou o surgimento do \textit{Model Context Protocol} (MCP). Este consolidou-se como o principal padrão de mercado para a integração entre agentes e ferramentas. O protocolo atingiu a marca de 97 milhões de \textit{downloads} mensais de seus Kits de Desenvolvimento de Software (SDK, do inglês \textit{Software Development Kit}) no início de 2026. Seu diferencial reside no uso de ferramentas (\textit{tools}) descritas em linguagem natural para a execução de ações dinâmicas. Contudo, essa flexibilidade introduz uma dependência crítica da qualidade documental. Assim como \textit{code smells} comprometem a manutenibilidade de \textit{softwares} tradicionais, descrições ambíguas no contexto do MCP podem induzir agentes a erros e alucinações. A alucinação caracteriza-se pela fabricação de comandos, referências ou conteúdos inexistentes, substituindo fatos reais por informações imaginadas extraídas da base de treinamento do modelo.

Estudos recentes investigam a prevalência de anomalias (ou \textit{smells}) nas descrições de ferramentas MCP por meio da \textbf{análise textual}. Esta técnica avalia se a documentação comunica claramente o propósito, parâmetros e restrições da ferramenta, verificando sua clareza e completude sem analisar a implementação técnica \cite{mcpAtFirstGlance}. Entretanto, evidências indicam que falhas de configuração e documentação somam 38,66\% dos problemas em servidores MCP \cite{Taraghi2026}. Se a precisão desses textos depende de processos manuais desacoplados do código, surge o risco de discrepância semântica. Diante disso, o \textbf{problema central abordado neste estudo reside na carência de investigações empíricas que confrontem sistematicamente as descrições com a implementação real no código-fonte, a fim de mitigar divergências técnicas que comprometem a confiabilidade de sistemas agênticos}.

A análise desse problema é fundamental para a evolução sustentável de sistemas agênticos. Inconsistências estruturais reduzem a eficácia dos agentes de IA e expõem o sistema a vulnerabilidades, como a injeção de comandos \cite{Taraghi2026}. O ecossistema MCP enfrenta desafios de manutenção e segurança. Investigações empíricas identificam a prevalência de anomalias em 97,1\% das ferramentas disponíveis \cite{mcpAtFirstGlance}. Relatórios de segurança institucional alertam que ambiguidades na documentação permitem a execução de ações não autorizadas ou a exposição de dados sensíveis \cite{NSA2026}. Diante desses riscos, a avaliação de \textit{smells} enriquecida pelo código-fonte resolve a lacuna de confiabilidade ao confrontar a descrição textual com a implementação real. Essa validação técnica direta mitiga a divida técnica e assegura o desenvolvimento de integrações mais robustas.

Para cobrir a lacuna de fidelidade entre a documentação e a implementação real no ecossistema agêntico, este trabalho tem como \textbf{objetivo geral realizar uma avaliação empírica comparativa da qualidade de descrições de ferramentas MCP, integrando a análise textual clássica ao contexto estrutural do código-fonte}. Busca-se validar a eficacia da \textbf{avaliação enriquecida} na detecção de discrepâncias semânticas e \textit{smells} para apoiar a confiabilidade de integrações baseadas em modelos de linguagem. Para o alcance dessa meta, definem-se os seguintes objetivos específicos: (i) mensurar o nível de clareza e completude das descrições em linguagem natural a partir de rubricas multidimensionais de qualidade fundamentadas na literatura base \cite{2026arXiv260214878M}; (ii) quantificar o impacto da incorporação de propriedades estruturais, extraídas via grafo de chamadas, na identificação de omissões e inconsistências técnicas não detectáveis exclusivamente por análise textual; (iii) determinar o ganho de precisão diagnóstica proporcionado pelo contexto técnico, fornecendo subsídio estatísticos para a mitigação da dívida técnica e para o fortalecimento da robustez em sistemas baseados em MCP.

Como resultados, espera-se que a avaliação enriquecida identifique inconsistências ocultas em 206 servidores MCP minerados do GitHub que possuem mais de 100 estrelas. A adoção desse limiar de popularidade assegura a análise de projetos maduros e mitiga o ruído de repositórios experimentais \cite{Borges2016, 2026arXiv260214878M}. A investigação focará no confronto direto entre as descrições em linguagem natural e a lógica real de execução das ferramentas, extraída por meio de grafos de chamadas (\textit{call graphs}) com profundidade de três níveis. Espera-se revelar discrepâncias técnicas, como parâmetros existentes no código mas ausentes na documentação, ou funcionalidades descritas que não possuem correspondência na implementação. Assim, busca-se obter evidências que validem a necessidade de integrar o código-fonte ao processo de análise, demonstrando que métricas exclusivamente textuais são insuficientes para a identificação precisa de \textit{smells} que comprometem a confiabilidade do sistema \cite{Taraghi2026}.

As próximas seções deste artigo seguem a seguinte organização: a Seção 2 apresenta a Fundamentação Teórica sobre MCP e detecção de \textit{smells}. A Seção 3 discute os trabalhos relacionados. Por fim, a Seção 4 descreve os materiais e métodos empregados na replicação e extensão do experimento.
